\documentclass[11pt,a4paper]{article}

\usepackage[T1]{fontenc}
\usepackage{libertinus}
\usepackage{microtype}
\usepackage[a4paper,margin=26mm,headheight=23pt]{geometry}
\usepackage{amsmath,amssymb,mathtools,amsthm}
\usepackage{mathrsfs}
\usepackage{bm}
\usepackage{booktabs,longtable,tabularx,array,multirow}
\usepackage{enumitem}
\usepackage{xcolor}
\usepackage[most]{tcolorbox}
\usepackage{listings}
\usepackage{fancyhdr}
\usepackage{hyperref}
\usepackage[nameinlink,capitalise,noabbrev]{cleveref}
\usepackage{caption}
\usepackage{needspace}
\usepackage{ragged2e}
\usepackage{url}

\definecolor{Ink}{HTML}{17212B}
\definecolor{Blue}{HTML}{245B8A}
\definecolor{Teal}{HTML}{207A74}
\definecolor{Gold}{HTML}{A66A00}
\definecolor{Red}{HTML}{9A3D3D}
\definecolor{PaleBlue}{HTML}{EFF6FB}
\definecolor{PaleTeal}{HTML}{EEF8F6}
\definecolor{PaleGold}{HTML}{FFF8E8}
\definecolor{PaleRed}{HTML}{FBEFEF}
\definecolor{Rule}{HTML}{CBD5DF}

\hypersetup{
  colorlinks=true,
  linkcolor=Blue,
  citecolor=Teal,
  urlcolor=Blue,
  pdftitle={All-Orders Asymptotic Transseries for the Swing Factorial and Its Inverse},
  pdfauthor={OpenAI},
  pdfsubject={Swing factorial, gamma ratios, Bernoulli and Bell polynomials, Lambert W, asymptotic inversion}
}

\pagestyle{fancy}
\fancyhf{}
\fancyhead[L]{\small\textcolor{Ink}{Swing factorial transseries}}
\fancyhead[R]{\small\textcolor{Ink}{\nouppercase{\leftmark}}}
\fancyfoot[C]{\small\thepage}
\renewcommand{\headrulewidth}{0.4pt}
\renewcommand{\headrule}{\hbox to\headwidth{\color{Rule}\leaders\hrule height \headrulewidth\hfill}}

\setcounter{tocdepth}{2}
\setcounter{secnumdepth}{3}
\setlength{\parindent}{0pt}
\setlength{\parskip}{0.55em}
\setlist{nosep,leftmargin=2em}
\allowdisplaybreaks

\newtheoremstyle{clean}%
  {0.7em}{0.7em}{\itshape}{}% space above/below, body, indent
  {\color{Blue}\bfseries}{.}{0.55em}{}
\theoremstyle{clean}
\newtheorem{theorem}{Theorem}[section]
\newtheorem{proposition}[theorem]{Proposition}
\newtheorem{lemma}[theorem]{Lemma}
\newtheorem{corollary}[theorem]{Corollary}
\newtheorem{conjecture}[theorem]{Conjecture}
\theoremstyle{definition}
\newtheorem{definition}[theorem]{Definition}
\newtheorem{remark}[theorem]{Remark}
\newtheorem{example}[theorem]{Example}

\newtcolorbox{summarybox}[1][]{
  enhanced,breakable,colback=PaleBlue,colframe=Blue,
  boxrule=0.7pt,arc=1.5mm,left=2mm,right=2mm,top=1.5mm,bottom=1.5mm,
  title={#1},fonttitle=\bfseries
}
\newtcolorbox{resultbox}[1][]{
  enhanced,breakable,colback=PaleTeal,colframe=Teal,
  boxrule=0.8pt,arc=1.5mm,left=2mm,right=2mm,top=1.5mm,bottom=1.5mm,
  title={#1},fonttitle=\bfseries
}
\newtcolorbox{warningbox}[1][]{
  enhanced,breakable,colback=PaleGold,colframe=Gold,
  boxrule=0.7pt,arc=1.5mm,left=2mm,right=2mm,top=1.5mm,bottom=1.5mm,
  title={#1},fonttitle=\bfseries
}
\newtcolorbox{auditbox}[1][]{
  enhanced,breakable,colback=PaleRed,colframe=Red,
  boxrule=0.7pt,arc=1.5mm,left=2mm,right=2mm,top=1.5mm,bottom=1.5mm,
  title={#1},fonttitle=\bfseries
}

\newcommand{\sw}{\operatorname{sw}}
\newcommand{\e}{\mathrm e}
\newcommand{\ii}{\mathrm i}
\newcommand{\eps}{\varepsilon}
\newcommand{\al}{\alpha}
\newcommand{\be}{\beta}
\newcommand{\B}{\mathrm B}
\newcommand{\W}{\mathrm W}
\newcommand{\Coeff}[2]{\left[#1\right]#2}
\newcommand{\OO}{\mathcal O}
\newcommand{\R}{\mathbb R}
\newcommand{\C}{\mathbb C}
\newcommand{\N}{\mathbb N}
\newcommand{\Dop}{\mathscr D}
\newcommand{\calH}{\mathcal H}
\newcommand{\calS}{\mathcal S}
\newcommand{\calM}{\mathcal M}
\newcommand{\calP}{\mathcal P}
\newcommand{\sgn}{\operatorname{sgn}}
\newcommand{\defeq}{\mathrel{:=}}
\newcommand{\given}{\,\middle|\,}
\DeclareMathOperator{\Li}{Li}
\DeclareMathOperator{\argmin}{arg\,min}

\lstdefinestyle{code}{
  basicstyle=\ttfamily\small,
  backgroundcolor=\color{PaleBlue},
  frame=single,
  rulecolor=\color{Rule},
  breaklines=true,
  columns=fullflexible,
  keepspaces=true,
  showstringspaces=false,
  xleftmargin=1em,
  xrightmargin=1em,
  aboveskip=0.8em,
  belowskip=0.8em
}

\title{\vspace{-1.5cm}\textbf{All-Orders Asymptotic Transseries for the Swing Factorial and Its Two Inverse Branches}\\[0.35em]
\large Bernoulli--Bell coefficient calculus, Lambert--\(W\) normal forms,\\[-0.05em]
\large logarithmic reversion, and beyond-all-orders control}
\author{OpenAI\\\small A standalone research note prepared in the notation of the ProveIt\\[-0.1em]
\small coefficient-calculus and transseries drafts}
\date{3 September 2026}

\begin{document}
\maketitle

\begin{abstract}
The swing factorial, OEIS A056040,
\[
  \sw(n)=\frac{n!}{\lfloor n/2\rfloor!^{\,2}},
\]
is not asymptotically represented by one smooth branch: its even and odd subsequences have different power prefactors and the full sequence repeatedly rises and falls.  We therefore split it into the two gamma-ratio interpolants
\[
  S_\eps(m)=\frac{\Gamma(2m+\eps+1)}{\Gamma(m+1)^2},\qquad \eps\in\{0,1\},
\]
so that \(S_\eps(m)=\sw(2m+\eps)\) for integral \(m\geq0\).  For each branch we derive the complete Poincar\'e transseries
\[
 S_\eps(m)\sim K_\eps 4^m m^{\be_\eps}
 \exp\!\left(\sum_{r\geq1}\lambda_r^{(\eps)}m^{-r}\right)
 =K_\eps4^m m^{\be_\eps}\sum_{r\geq0}c_r^{(\eps)}m^{-r},
\]
where \(K_\eps=2^\eps/\sqrt\pi\), \(\be_\eps=\eps-\tfrac12\), the logarithmic coefficients are explicit Bernoulli-polynomial differences, and the multiplicative coefficients are complete Bell polynomials.  We give partition sums, triangular recurrences, parity-transfer identities, exact coefficient tables, rigorous fixed-order remainders, and an exact Binet-integral completion.

Each \(S_\eps\) is strictly increasing on the relevant half-line and hence has a real inverse \(M_\eps\).  Its dominant power--exponential phase is inverted exactly by a Lambert function: \(\W_{-1}\) for the even sheet and \(\W_0\) for the odd sheet.  Writing that Lambert core as \(w_\eps(y)\), we prove the all-orders inverse expansion
\[
  M_\eps(y)\sim w_\eps(y)+\sum_{r\geq1}d_r^{(\eps)}w_\eps(y)^{-r}.
\]
The coefficients \(d_r^{(\eps)}\) are specified in three equivalent ways: a triangular coefficient-extraction recurrence, a closed Lagrange--B\"urmann differential formula, and a finite coefficient-array algorithm.  Finally, we unfold the Lambert core into the ordinary polynomial--logarithmic scale
\[
 M_\eps(y)\sim \frac1{\log4}\left(X_\eps-\be_\eps\log X_\eps+
 \sum_{r\geq1}\frac{P_r^{(\eps)}(\log X_\eps)}{X_\eps^r}\right),
\]
with an all-orders recurrence for every coefficient of every polynomial \(P_r^{(\eps)}\).  Numerical residual tests confirm the formulae to the displayed orders.
\end{abstract}

\begin{summarybox}[Interpretation of ``full transseries'']
The word \emph{full} is used in three precise senses here: all algebraic orders are given by explicit coefficient formulae; both parity sheets are included; and the exact Binet integral is retained as the canonical beyond-all-orders completion.  A further decomposition of the optimally truncated Binet remainder into terminant/Stokes sectors is representation-dependent.  The transport rule for any such exponentially improved sector is nevertheless derived in \cref{sec:beyond}.
\end{summarybox}

\tableofcontents
\newpage

\section{Introduction and statement of the main results}
\label{sec:intro}

The swing factorial appears deceptively close to the central binomial coefficient, yet its parity structure changes the inverse problem qualitatively.  At even arguments it \emph{is} the central binomial coefficient; at odd arguments it is that coefficient multiplied by a linear factor.  Thus the direct asymptotics require a period-two transseries, while an inverse of the entire sequence does not exist in the usual single-valued sense.

The derivation has four layers.
\begin{enumerate}
\item The duplication formula reduces both parity branches to a common gamma ratio.
\item Shifted Stirling expansions express the logarithmic coefficients as Bernoulli-polynomial differences.  Exponentiation is then handled exactly by complete Bell polynomials.
\item The leading inverse phase \(\al m+\be\log m\) is solved exactly with a branch of Lambert \(W\).
\item Lagrange--B\"urmann inversion around that exact carrier produces the complete inverse-power correction grid.  A second reversion unfolds \(W\) into powers of \(1/\log y\) with polynomial coefficients in \(\log\log y\).
\end{enumerate}

The two carrier constants used throughout are
\begin{equation}
  \al=\log4,\qquad
  K_\eps=\frac{2^\eps}{\sqrt\pi},\qquad
  \be_\eps=\eps-\frac12,
  \qquad \eps\in\{0,1\}.
  \label{eq:carrier-constants}
\end{equation}

\begin{resultbox}[Main direct and inverse formulae]
For \(m\to+\infty\),
\begin{align}
 S_\eps(m)
 &\sim K_\eps\e^{\al m}m^{\be_\eps}
 \exp\!\left(\sum_{r\ge1}\lambda_r^{(\eps)}m^{-r}\right),
 \label{eq:summary-direct-log}\\
 \lambda_r^{(\eps)}
 &=\frac{(-1)^{r+1}}{r(r+1)}
 \left\{\B_{r+1}\!\left(\eps+\frac12\right)-\B_{r+1}(1)\right\}.
 \label{eq:summary-lambda}
\end{align}
Writing \(\exp(\sum_{r\ge1}\lambda_rt^r)=\sum_{r\ge0}c_rt^r\),
\begin{equation}
 c_r^{(\eps)}=\frac1{r!}
 Y_r\!\left(1!\lambda_1^{(\eps)},2!\lambda_2^{(\eps)},\ldots,r!\lambda_r^{(\eps)}\right),
 \label{eq:summary-c}
\end{equation}
where \(Y_r\) is the complete exponential Bell polynomial.

Let \(M_\eps=S_\eps^{-1}\), and set
\begin{align}
 w_0(y)&=-\frac1{2\al}\,
 \W_{-1}\!\left(-\frac{2\al}{\pi y^2}\right),
 \label{eq:summary-w0}\\
 w_1(y)&=\frac1{2\al}\,
 \W_0\!\left(\frac{\pi\al}{2}y^2\right).
 \label{eq:summary-w1}
\end{align}
Then
\begin{equation}
 M_\eps(y)\sim w_\eps(y)+\sum_{r\ge1}d_r^{(\eps)}w_\eps(y)^{-r}.
 \label{eq:summary-inverse}
\end{equation}
If \(H_\eps(t)=\sum_{r\ge1}\lambda_r^{(\eps)}t^r\),
\(P_\eps(t)=\al+\be_\eps t\), and
\begin{equation}
 \Dop_\eps=-\frac{t^2}{P_\eps(t)}\frac{d}{dt},
 \label{eq:summary-Dop}
\end{equation}
then the closed coefficient formula is
\begin{equation}
 d_r^{(\eps)}=
 \Coeff{t^r}{
 \sum_{q=1}^{\lfloor(r+1)/2\rfloor}
 \frac{(-1)^q}{q!}\,
 \Dop_\eps^{\,q-1}
 \left(\frac{H_\eps(t)^q}{P_\eps(t)}\right)}.
 \label{eq:summary-d}
\end{equation}
The inverse in the original index is \(N_\eps(y)=2M_\eps(y)+\eps\).
\end{resultbox}

\section{Exact structure and the unavoidable parity split}
\label{sec:exact}

\subsection{Equivalent exact definitions}

For \(n\in\N_0\), define
\begin{equation}
 \sw(n)=\frac{n!}{\lfloor n/2\rfloor!^{\,2}}.
 \label{eq:swing-def}
\end{equation}
This is OEIS A056040.  Separating parity gives the elementary identities
\begin{align}
 \sw(2m)&=\binom{2m}{m},
 \label{eq:even-exact}\\
 \sw(2m+1)&=(2m+1)\binom{2m}{m}.
 \label{eq:odd-exact}
\end{align}
Equivalently,
\begin{equation}
 \sw(n)=2^{n-(n\bmod2)}\prod_{k=1}^n k^{(-1)^{k+1}}.
 \label{eq:swing-product}
\end{equation}
The first terms are
\[
 1,1,2,6,6,30,20,140,70,630,252,2772,\ldots.
\]

For later analytic work it is convenient to interpolate the two subsequences separately.

\begin{definition}[Parity branches]
For \(m\ge0\) and \(\eps\in\{0,1\}\), set
\begin{equation}
 S_\eps(m)=\frac{\Gamma(2m+\eps+1)}{\Gamma(m+1)^2}.
 \label{eq:S-def}
\end{equation}
Then \(S_\eps(m)=\sw(2m+\eps)\) whenever \(m\) is a nonnegative integer.
\end{definition}

The gamma duplication formula gives the common normal form
\begin{equation}
 S_\eps(m)=K_\eps4^m
 \frac{\Gamma(m+\eps+\tfrac12)}{\Gamma(m+1)},
 \qquad K_\eps=\frac{2^\eps}{\sqrt\pi}.
 \label{eq:duplication-normal-form}
\end{equation}
In particular,
\begin{equation}
 S_1(m)=(2m+1)S_0(m).
 \label{eq:branch-transfer-exact}
\end{equation}

\subsection{Generating functions as a consistency check}

The exponential generating function follows immediately from \cref{eq:even-exact,eq:odd-exact}:
\begin{align}
 \sum_{n\ge0}\sw(n)\frac{z^n}{n!}
 &=\sum_{m\ge0}\frac{z^{2m}}{(m!)^2}
   +z\sum_{m\ge0}\frac{z^{2m}}{(m!)^2} \\
 &=(1+z)I_0(2z),
 \label{eq:egf}
\end{align}
where \(I_0\) is the modified Bessel function.  The ordinary generating function is
\begin{equation}
 \sum_{n\ge0}\sw(n)z^n
 =\frac1{\sqrt{1-4z^2}}+\frac{z}{(1-4z^2)^{3/2}}.
 \label{eq:ogf}
\end{equation}
These exact formulae are not needed for the gamma-ratio derivation, but they independently display the two singular amplitudes that later become the powers \(m^{-1/2}\) and \(m^{1/2}\).

\subsection{Why there is no one-valued inverse of the sequence}

The adjacent ratios are exact:
\begin{align}
 \frac{\sw(2m+1)}{\sw(2m)}&=2m+1,
 \label{eq:up-ratio}\\
 \frac{\sw(2m+2)}{\sw(2m+1)}&=\frac{2}{m+1}.
 \label{eq:down-ratio}
\end{align}
Thus the sequence rises steeply from every even term to the following odd term and, for \(m\ge2\), falls from that odd term to the next even term.  Consequently there is no eventual monotonicity and no ordinary global inverse \(n=n(y)\).

\begin{warningbox}[The inverse is intrinsically two-sheeted]
Any formula advertised as \emph{the} inverse swing factorial without specifying parity, a smoothing convention, or a generalized-inverse convention is incomplete.  The canonical analytic object is the pair
\[
 M_0=S_0^{-1},\qquad M_1=S_1^{-1},
\]
or, in the original index variable,
\[
 N_0(y)=2M_0(y),\qquad N_1(y)=2M_1(y)+1.
\]
\end{warningbox}

\subsection{Monotonicity of each analytic sheet}

\begin{proposition}\label{prop:monotonicity}
The function \(S_0\) is strictly increasing on \((0,\infty)\), with \(S_0(0)=1\), and \(S_1\) is strictly increasing on \([0,\infty)\), with \(S_1(0)=1\).  Hence each branch maps \([0,\infty)\) continuously onto \([1,\infty)\) and has a unique real inverse \(M_\eps:[1,\infty)\to[0,\infty)\).
\end{proposition}

\begin{proof}
Differentiating \cref{eq:S-def} logarithmically gives
\begin{equation}
 \frac{S_\eps'(m)}{S_\eps(m)}
 =2\psi(2m+\eps+1)-2\psi(m+1),
 \label{eq:log-derivative}
\end{equation}
where \(\psi=\Gamma'/\Gamma\).  For \(x>0\),
\(
 \psi'(x)=\sum_{k\ge0}(x+k)^{-2}>0,
\)
so \(\psi\) is strictly increasing.  If \(\eps=0\), the two arguments in \cref{eq:log-derivative} coincide only at \(m=0\) and the first is larger for \(m>0\).  If \(\eps=1\), the first is larger for every \(m\ge0\).  The range assertion follows from continuity, the initial values, and the exponential growth established below.
\end{proof}

\section{The coefficient-calculus toolkit}
\label{sec:toolkit}

\subsection{Bernoulli polynomials}

We use the convention
\begin{equation}
 \frac{te^{xt}}{e^t-1}=\sum_{n\ge0}\B_n(x)\frac{t^n}{n!},
 \qquad \B_1=-\frac12.
 \label{eq:Bernoulli-egf}
\end{equation}
The two identities needed repeatedly are
\begin{equation}
 \B_n\!\left(\frac12\right)=\bigl(2^{1-n}-1\bigr)\B_n,
 \qquad
 \B_n(1)=\B_n\quad(n\ne1).
 \label{eq:Bernoulli-half}
\end{equation}

\subsection{Complete Bell polynomials}

The complete exponential Bell polynomials \(Y_n\) are defined by
\begin{equation}
 \exp\!\left(\sum_{k\ge1}x_k\frac{t^k}{k!}\right)
 =\sum_{n\ge0}Y_n(x_1,\ldots,x_n)\frac{t^n}{n!}.
 \label{eq:Bell-def}
\end{equation}
Consequently, if
\(
 H(t)=\sum_{k\ge1}\lambda_kt^k
\)
and
\(
 C(t)=e^{H(t)}=\sum_{n\ge0}c_nt^n,
\)
then
\begin{align}
 c_n&=\frac1{n!}Y_n(1!\lambda_1,2!\lambda_2,\ldots,n!\lambda_n),
 \label{eq:Bell-c}\\
 &=\sum_{\substack{j_1,j_2,\ldots\ge0\\
                    j_1+2j_2+\cdots=n}}
   \prod_{k=1}^n\frac{\lambda_k^{j_k}}{j_k!},
 \label{eq:partition-c}\\
 c_0&=1,\qquad
 c_n=\frac1n\sum_{k=1}^n k\lambda_kc_{n-k}.
 \label{eq:recurrence-c}
\end{align}
The first formula is the compact Bell-polynomial form, the second is the frequency-vector or partition form, and the third is the computational recurrence obtained from \(C'=H'C\).

\subsection{Formal and asymptotic meanings}

An identity between coefficient series in the variable \(t\) is first interpreted formally.  When \(t=m^{-1}\) with \(m\to+\infty\), the same algebra gives a Poincar\'e asymptotic expansion: truncating after \(t^{N-1}\) leaves a remainder \(\OO(t^N)\).  This distinction matters because the Bernoulli-generated series below is generally divergent.  Formal coefficient identities remain exact even when analytic convergence fails.

\begin{auditbox}[Normalization audit]
The symbols \(\B_n(x)\), \(Y_n\), and \(\W_k\) denote Bernoulli polynomials, complete exponential Bell polynomials, and Lambert-function branches, respectively.  Keeping these three unrelated uses of ``B/W-type'' notation separate prevents the factorial and sign errors that often occur in asymptotic exponentiation and inversion.
\end{auditbox}

\section{All-orders direct transseries}
\label{sec:direct}

\subsection{Logarithmic expansion}

The shifted Stirling expansion states, for fixed \(h\),
\begin{equation}
 \log\Gamma(m+h)
 \sim \left(m+h-\frac12\right)\log m-m+\frac12\log(2\pi)
 +\sum_{k\ge2}\frac{(-1)^k\B_k(h)}{k(k-1)m^{k-1}}.
 \label{eq:shifted-Stirling}
\end{equation}
Subtracting the instances \(h=\eps+\tfrac12\) and \(h=1\) in
\cref{eq:duplication-normal-form} gives the central direct theorem.

\begin{theorem}[Logarithmic swing-factorial transseries]
\label{thm:log-direct}
For each \(\eps\in\{0,1\}\),
\begin{equation}
 \log S_\eps(m)
 \sim \log K_\eps+\al m+\be_\eps\log m
 +\sum_{r\ge1}\frac{\lambda_r^{(\eps)}}{m^r},
 \label{eq:log-direct}
\end{equation}
where
\begin{equation}
 \boxed{\displaystyle
 \lambda_r^{(\eps)}=
 \frac{(-1)^{r+1}}{r(r+1)}
 \left[
 \B_{r+1}\!\left(\eps+\frac12\right)-\B_{r+1}(1)
 \right].}
 \label{eq:lambda-general}
\end{equation}
For every fixed \(N\ge1\), truncation through \(r=N-1\) has error \(\OO(m^{-N})\) as \(m\to+\infty\).  The same assertion holds uniformly in every closed subsector \(|\arg m|\le\pi-\delta\) that stays away from the gamma poles and chosen logarithmic cuts.
\end{theorem}

\begin{proof}
Insert \cref{eq:shifted-Stirling} with \(h=\eps+\tfrac12\) and \(h=1\) into
\begin{equation*}
 \log S_\eps(m)=\log K_\eps+\al m+
 \log\Gamma\!\left(m+\eps+\frac12\right)-\log\Gamma(m+1).
\end{equation*}
The terms \(-m\) and \(\tfrac12\log(2\pi)\) cancel.  The coefficient of \(\log m\) is
\((\eps+\tfrac12-\tfrac12)-(1-\tfrac12)=\eps-\tfrac12=\be_\eps\).
Writing \(r=k-1\) in the remaining Bernoulli sum gives \cref{eq:lambda-general}.  The remainder statement is the corresponding fixed-order remainder theorem for shifted Stirling expansions, applied to a fixed difference of two shifts.
\end{proof}

For the even sheet, \cref{eq:Bernoulli-half} simplifies the answer considerably:
\begin{equation}
 \lambda_r^{(0)}=
 \frac{(-1)^{r+1}(2^{-r}-2)\B_{r+1}}{r(r+1)}.
 \label{eq:lambda-even}
\end{equation}
Hence
\begin{equation}
 \lambda_{2j}^{(0)}=0\qquad(j\ge1).
 \label{eq:lambda-even-zero}
\end{equation}
The exact branch-transfer identity \cref{eq:branch-transfer-exact} yields a second useful simplification:
\begin{align}
 \log S_1(m)-\log S_0(m)
 &=\log(2m+1)=\log2+\log m+\log\!\left(1+\frac1{2m}\right),
 \nonumber\\[-0.2em]
 \lambda_r^{(1)}
 &=\lambda_r^{(0)}+\frac{(-1)^{r+1}}{r2^r}.
 \label{eq:lambda-transfer}
\end{align}

The first logarithmic coefficients are displayed in \cref{tab:lambda}.

\begin{table}[htbp]
\centering
\caption{Logarithmic coefficients \(\lambda_r^{(\eps)}\).}
\label{tab:lambda}
\begin{tabular}{@{}rcc@{}}
\toprule
\(r\) & \(\lambda_r^{(0)}\) & \(\lambda_r^{(1)}\)\\
\midrule
1  & \(-1/8\)       & \(3/8\)\\
2  & \(0\)          & \(-1/8\)\\
3  & \(1/192\)      & \(3/64\)\\
4  & \(0\)          & \(-1/64\)\\
5  & \(-1/640\)     & \(3/640\)\\
6  & \(0\)          & \(-1/384\)\\
7  & \(17/14336\)   & \(33/14336\)\\
8  & \(0\)          & \(-1/2048\)\\
9  & \(-31/18432\)  & \(-3/2048\)\\
10 & \(0\)          & \(-1/10240\)\\
\bottomrule
\end{tabular}
\end{table}

\subsection{Multiplicative expansion and general coefficient formulae}

Define
\begin{equation}
 H_\eps(t)=\sum_{r\ge1}\lambda_r^{(\eps)}t^r,
 \qquad
 C_\eps(t)=\exp(H_\eps(t))=\sum_{r\ge0}c_r^{(\eps)}t^r.
 \label{eq:H-C-def}
\end{equation}

\begin{theorem}[Multiplicative swing-factorial transseries]
\label{thm:mult-direct}
For \(\eps\in\{0,1\}\),
\begin{equation}
 S_\eps(m)\sim K_\eps4^m m^{\be_\eps}
 \sum_{r\ge0}\frac{c_r^{(\eps)}}{m^r},
 \label{eq:mult-direct}
\end{equation}
where any of the equivalent all-orders formulae
\begin{align}
 c_r^{(\eps)}
 &=\frac1{r!}Y_r\!\left(
 1!\lambda_1^{(\eps)},2!\lambda_2^{(\eps)},\ldots,r!\lambda_r^{(\eps)}
 \right),
 \label{eq:c-Bell}\\
 &=\sum_{\substack{j_1,j_2,\ldots\ge0\\j_1+2j_2+\cdots=r}}
 \prod_{k=1}^r\frac{(\lambda_k^{(\eps)})^{j_k}}{j_k!},
 \label{eq:c-partition}\\
 c_0^{(\eps)}&=1,\qquad
 c_r^{(\eps)}=\frac1r\sum_{k=1}^r
 k\lambda_k^{(\eps)}c_{r-k}^{(\eps)}
 \label{eq:c-recurrence}
\end{align}
compute every coefficient.  For fixed \(N\),
\begin{equation}
 S_\eps(m)=K_\eps4^m m^{\be_\eps}
 \left(\sum_{r=0}^{N-1}c_r^{(\eps)}m^{-r}+\OO(m^{-N})\right).
 \label{eq:mult-remainder}
\end{equation}
\end{theorem}

\begin{proof}
Exponentiate \cref{eq:log-direct} formally and apply \cref{eq:Bell-c,eq:partition-c,eq:recurrence-c}.  For a fixed truncation order, write the logarithmic remainder as \(R_N(m)=\OO(m^{-N})\).  Then \(e^{R_N(m)}=1+\OO(m^{-N})\), so formal exponentiation also gives the claimed analytic remainder after collecting terms through order \(m^{-(N-1)}\).
\end{proof}

Because \(S_1=(2m+1)S_0=2m(1+(2m)^{-1})S_0\), the multiplicative coefficients satisfy the particularly simple transfer law
\begin{equation}
 \boxed{\displaystyle
 c_r^{(1)}=c_r^{(0)}+\frac12c_{r-1}^{(0)},
 \qquad c_{-1}^{(0)}=0.}
 \label{eq:c-transfer}
\end{equation}
This identity is exact at the level of formal asymptotic series and is a useful independent check on every computed coefficient.

The first terms are
\begin{align}
 S_0(m)&\sim \frac{4^m}{\sqrt{\pi m}}
 \left(1-\frac1{8m}+\frac1{128m^2}+\frac5{1024m^3}\right.
 \nonumber\\[-0.25em]
 &\hspace{34mm}\left.
 -\frac{21}{32768m^4}-\frac{399}{262144m^5}
 +\frac{869}{4194304m^6}+\cdots\right),
 \label{eq:even-direct-first}\\
 S_1(m)&\sim 2\,4^m\sqrt{\frac{m}{\pi}}
 \left(1+\frac3{8m}-\frac7{128m^2}+\frac9{1024m^3}\right.
 \nonumber\\[-0.25em]
 &\hspace{34mm}\left.
 +\frac{59}{32768m^4}-\frac{483}{262144m^5}
 -\frac{2323}{4194304m^6}+\cdots\right).
 \label{eq:odd-direct-first}
\end{align}

\begin{table}[htbp]
\centering
\caption{Multiplicative coefficients \(c_r^{(\eps)}\).}
\label{tab:c}
\small
\begin{tabular}{@{}rcc@{}}
\toprule
\(r\) & \(c_r^{(0)}\) & \(c_r^{(1)}\)\\
\midrule
0 & \(1\) & \(1\)\\
1 & \(-1/8\) & \(3/8\)\\
2 & \(1/128\) & \(-7/128\)\\
3 & \(5/1024\) & \(9/1024\)\\
4 & \(-21/32768\) & \(59/32768\)\\
5 & \(-399/262144\) & \(-483/262144\)\\
6 & \(869/4194304\) & \(-2323/4194304\)\\
7 & \(39325/33554432\) & \(42801/33554432\)\\
8 & \(-334477/2147483648\) & \(923923/2147483648\)\\
9 & \(-28717403/17179869184\) & \(-30055311/17179869184\)\\
10 & \(59697183/274877906944\) & \(-170042041/274877906944\)\\
\bottomrule
\end{tabular}
\end{table}

\subsection{A single period-two formula in the original index}

For an integer \(n\), let
\begin{equation}
 \chi_\eps(n)=\frac{1+(-1)^{n-\eps}}2,
 \qquad \eps\in\{0,1\}.
 \label{eq:parity-projector}
\end{equation}
Then \(\chi_\eps(n)\) is the indicator of \(n\equiv\eps\pmod2\).  Since \(m=(n-\eps)/2\) on that parity class, \cref{eq:mult-direct} becomes the unified period-two transseries
\begin{equation}
 \boxed{\displaystyle
 \sw(n)\sim
 \sum_{\eps=0}^1\chi_\eps(n)\,
 \frac{2^n}{\sqrt\pi}
 \left(\frac{n-\eps}{2}\right)^{\eps-1/2}
 \sum_{r\ge0}c_r^{(\eps)}
 \left(\frac{n-\eps}{2}\right)^{-r}.}
 \label{eq:unified-n-direct}
\end{equation}
The projectors are not cosmetic: they encode two genuinely different power laws multiplying the common exponential carrier \(2^n\).

\subsection{Late terms, divergence, and the natural exponential scale}
\label{subsec:divergence}

The Bernoulli numbers satisfy
\(
 \B_{2j}=(-1)^{j-1}2(2j)!\zeta(2j)/(2\pi)^{2j}
\).
For the nonzero even-sheet logarithmic coefficients this gives
\begin{equation}
 \lambda_{2j-1}^{(0)}
 =\frac{(2^{1-2j}-2)\B_{2j}}{(2j-1)(2j)}
 \sim (-1)^j\frac{4(2j-2)!}{(2\pi)^{2j}}.
 \label{eq:lambda-late}
\end{equation}
The additional term in \cref{eq:lambda-transfer} is merely geometric and therefore does not alter the factorial late-term law on the odd sheet.  The logarithmic and multiplicative expansions are consequently divergent Gevrey-type asymptotic series.  Their least terms occur at an index proportional to \(2\pi m\), and the optimally truncated remainder has the familiar gamma-function scale associated with \(e^{-2\pi m}\) in magnitude on the positive real axis.  An exact integral that retains this information is given next.

\subsection{Exact Binet-integral completion}
\label{subsec:Binet}

For \(z>0\), Binet's second formula may be written
\begin{equation}
 \mathfrak L(z)=
 \left(z-\frac12\right)\log z-z+\frac12\log(2\pi)
 +2\int_0^\infty\frac{\arctan(t/z)}{e^{2\pi t}-1}\,dt
 =\log\Gamma(z).
 \label{eq:Binet}
\end{equation}
Therefore the exact, untruncated completion of \cref{eq:log-direct} is
\begin{equation}
 \boxed{\displaystyle
 \log S_\eps(m)=\log K_\eps+\al m+
 \mathfrak L\!\left(m+\eps+\frac12\right)-\mathfrak L(m+1).}
 \label{eq:exact-Binet-swing}
\end{equation}
Expanding the elementary part and the arctangent integral in inverse powers of \(m\) recovers \cref{eq:lambda-general}.  Retaining the integral instead of expanding it provides a canonical analytic remainder for numerical certification and for hyperasymptotic refinement.  In particular, no informal assignment of a Stokes constant is needed to make the real-axis formula exact.

\section{Lambert normal forms for the inverse branches}
\label{sec:Lambert}

\subsection{The dominant phase}

Let
\begin{equation}
 L_\eps(y)=\log\frac{y}{K_\eps},
 \qquad
 F_{0,\eps}(m)=\al m+\be_\eps\log m.
 \label{eq:L-F0}
\end{equation}
The dominant inverse problem is
\begin{equation}
 F_{0,\eps}(w)=L_\eps(y),
 \qquad\text{equivalently}\qquad
 K_\eps e^{\al w}w^{\be_\eps}=y.
 \label{eq:carrier-equation}
\end{equation}
For any nonzero \(\be\), exponentiating after division by \(\be\) gives
\begin{equation}
 w=\frac{\be}{\al}\,
 \W_k\!\left(
 \frac{\al}{\be}
 \left(\frac{y}{K}\right)^{1/\be}
 \right).
 \label{eq:generic-Lambert-core}
\end{equation}
The branch index \(k\) must be chosen so that \(w\to+\infty\).

\begin{proposition}[Real Lambert carriers]
\label{prop:Lambert-carriers}
For the even and odd swing-factorial sheets, the unique large positive solutions of \cref{eq:carrier-equation} are
\begin{align}
 w_0(y)&=-\frac1{2\al}\,
 \W_{-1}\!\left(-\frac{2\al}{\pi y^2}\right),
 \label{eq:w0}\\
 w_1(y)&=\frac1{2\al}\,
 \W_0\!\left(\frac{\pi\al}{2}y^2\right).
 \label{eq:w1}
\end{align}
Here \(w_0\) is real for sufficiently large \(y\), and \(w_1\) is real for every \(y>0\).
\end{proposition}

\begin{proof}
For \(\eps=0\), \(K_0=\pi^{-1/2}\) and \(\be_0=-1/2\).  Formula \cref{eq:generic-Lambert-core} becomes
\[
 w=-\frac1{2\al}\W_k\!\left(-\frac{2\al}{\pi y^2}\right).
\]
As \(y\to\infty\), the argument approaches \(0^-\).  The principal branch \(\W_0\) tends to zero and gives the small root of the approximate carrier; the branch \(\W_{-1}\) tends to \(-\infty\) and gives the required large positive root.  For \(\eps=1\), \(K_1=2/\sqrt\pi\) and \(\be_1=1/2\), which gives \cref{eq:w1}; its argument is positive, so the real branch is \(\W_0\).
\end{proof}

\begin{remark}[Why two different Lambert branches appear]
The even leading carrier \(e^{\al m}m^{-1/2}\) has a shallow minimum and two real inverse roots for sufficiently large levels if extended down to \(m>0\); the large root is encoded by \(\W_{-1}\).  The odd carrier \(e^{\al m}m^{1/2}\) is increasing and uses \(\W_0\).  This branch distinction is forced by the sign of \(\be_\eps\), not by a notational preference.
\end{remark}

\subsection{The displacement equation}

The full logarithmic equation is
\begin{equation}
 L_\eps(y)=F_{0,\eps}(m)+H_\eps(m^{-1}),
 \label{eq:full-log-inverse-equation}
\end{equation}
understood asymptotically to arbitrary algebraic order, or exactly after replacing the formal correction by the Binet expression in \cref{eq:exact-Binet-swing}.  Since \(w=w_\eps(y)\) solves \(F_{0,\eps}(w)=L_\eps(y)\), write
\begin{equation}
 m=w+D_\eps(t),\qquad t=w^{-1},
 \qquad D_\eps(t)=\sum_{r\ge1}d_r^{(\eps)}t^r.
 \label{eq:D-def}
\end{equation}
There is no constant displacement because the omitted logarithmic correction begins at order \(m^{-1}\).  Substitution into \cref{eq:full-log-inverse-equation} gives the fundamental formal identity
\begin{equation}
 \boxed{\displaystyle
 \al D_\eps(t)
 +\be_\eps\log\!\left(1+tD_\eps(t)\right)
 +H_\eps\!\left(\frac{t}{1+tD_\eps(t)}\right)=0.}
 \label{eq:fundamental-D}
\end{equation}
This one equation contains the entire inverse coefficient calculus.

\section{All-orders inverse coefficient calculus}
\label{sec:inverse-coefficients}

\subsection{Triangular recurrence}

For \(r\ge1\), let
\begin{equation}
 D_{\eps,<r}(t)=\sum_{j=1}^{r-1}d_j^{(\eps)}t^j.
 \label{eq:D-less-r}
\end{equation}
In the coefficient of \(t^r\) in the two nonlinear terms of \cref{eq:fundamental-D}, the unknown \(d_r^{(\eps)}\) cannot occur: it is always multiplied by at least one additional factor of \(t\).  Thus the recurrence is triangular.

\begin{theorem}[Triangular inverse recurrence]
\label{thm:d-recurrence}
The coefficients in \cref{eq:D-def} are uniquely determined by
\begin{equation}
 \boxed{\displaystyle
 d_r^{(\eps)}=-\frac1\al\Coeff{t^r}{
 \be_\eps\log\!\left(1+tD_{\eps,<r}(t)\right)
 +H_\eps\!\left(\frac{t}{1+tD_{\eps,<r}(t)}\right)}.}
 \label{eq:d-recurrence}
\end{equation}
Only \(\lambda_1^{(\eps)},\ldots,\lambda_r^{(\eps)}\) and the previously computed \(d_1^{(\eps)},\ldots,d_{r-1}^{(\eps)}\) are needed.
\end{theorem}

\begin{proof}
Take the coefficient of \(t^r\) in \cref{eq:fundamental-D}.  The linear term contributes \(\al d_r^{(\eps)}\).  Replacing \(D_\eps\) by \(D_{\eps,<r}\) in the remaining terms does not change their \(t^r\) coefficient, because the first occurrence of \(D\) is through \(tD\).  Solving the resulting linear equation proves the formula and uniqueness.
\end{proof}

\subsection{A closed Lagrange--B\"urmann differential formula}

The recurrence is computationally convenient, but a nonrecursive formula reveals the reversion mechanism.  Put
\begin{equation}
 P_\eps(t)=\al+\be_\eps t,
 \qquad
 \Dop_\eps=-\frac{t^2}{P_\eps(t)}\frac{d}{dt}.
 \label{eq:P-Dop}
\end{equation}
Since \(t=1/w\), the operator \(\Dop_\eps\) is precisely
\begin{equation}
 \Dop_\eps=\frac1{F_{0,\eps}'(w)}\frac{d}{dw}.
 \label{eq:Dop-w}
\end{equation}

\begin{lemma}[Perturbative Lagrange--B\"urmann formula]
\label{lem:perturbative-LB}
Let \(F_0'(w)\ne0\), and let \(x(q)\) solve
\[
 F_0(x(q))+q h(x(q))=F_0(w),\qquad x(0)=w.
\]
Then, as a formal series in \(q\),
\begin{equation}
 x(q)=w+\sum_{n\ge1}\frac{(-q)^n}{n!}
 \left(\frac1{F_0'(w)}\frac{d}{dw}\right)^{n-1}
 \left(\frac{h(w)^n}{F_0'(w)}\right).
 \label{eq:perturbative-LB}
\end{equation}
\end{lemma}

\begin{proof}
Use the phase coordinate \(u=F_0(x)\), write \(x=f(u)=F_0^{-1}(u)\), and set \(A(u)=h(f(u))\).  The equation becomes the near-identity reversion problem
\(u+qA(u)=F_0(w)\).  Lagrange--B\"urmann applied to the observable \(f\) gives
\[
 f(u)=f(R)+\sum_{n\ge1}\frac{(-q)^n}{n!}
 \left(\frac{d}{dR}\right)^{n-1}\!\left(f'(R)A(R)^n\right),
 \qquad R=F_0(w).
\]
Now \(f'(R)=1/F_0'(w)\) and \(d/dR=(1/F_0'(w))d/dw\), which is exactly \cref{eq:perturbative-LB}.
\end{proof}

\begin{theorem}[Closed formula for every inverse coefficient]
\label{thm:d-closed}
For \(r\ge1\),
\begin{equation}
 \boxed{\displaystyle
 d_r^{(\eps)}=
 \Coeff{t^r}{
 \sum_{q=1}^{\lfloor(r+1)/2\rfloor}
 \frac{(-1)^q}{q!}\,
 \Dop_\eps^{\,q-1}
 \left(\frac{H_\eps(t)^q}{P_\eps(t)}\right)}.}
 \label{eq:d-closed}
\end{equation}
The sum is finite for each \(r\).
\end{theorem}

\begin{proof}
Apply \cref{lem:perturbative-LB} with
\(
 F_0=F_{0,\eps}
\),
\(
 h(w)=H_\eps(w^{-1})
\), and then set \(q=1\).  Formula \cref{eq:Dop-w} converts the result into \cref{eq:d-closed}.  Since \(H_\eps(t)^q=\OO(t^q)\), while each application of \(\Dop_\eps\) raises the lowest power of \(t\) by one, the \(q\)-th summand begins at \(t^{2q-1}\).  Therefore it can contribute to \(d_r\) only if \(2q-1\le r\).
\end{proof}

\subsection{A completely finite coefficient-array algorithm}

Formula \cref{eq:d-closed} contains derivatives, but it can be reduced to finite sums of rational operations.  Define
\begin{equation}
 h_{q,s}^{(\eps)}=\Coeff{t^s}{H_\eps(t)^q}
 =q!\!\sum_{\substack{j_1,j_2,\ldots\ge0\\
              \sum_kj_k=q,\;\sum_kkj_k=s}}
 \prod_{k\ge1}\frac{(\lambda_k^{(\eps)})^{j_k}}{j_k!}.
 \label{eq:hqs}
\end{equation}
For fixed \(q\), introduce coefficients \(A_{q,j,r}^{(\eps)}\) by
\begin{equation}
 \Dop_\eps^{\,j}\!\left(\frac{H_\eps(t)^q}{P_\eps(t)}\right)
 =\sum_{r\ge0}A_{q,j,r}^{(\eps)}t^r.
 \label{eq:Aqjr-def}
\end{equation}
Since
\(
 P_\eps(t)^{-1}=\al^{-1}\sum_{v\ge0}(-\be_\eps/\al)^v t^v
\), the initialization and update are
\begin{align}
 A_{q,0,r}^{(\eps)}
 &=\frac1\al\sum_{s=q}^r
 \left(-\frac{\be_\eps}{\al}\right)^{r-s}h_{q,s}^{(\eps)},
 \label{eq:A-initial}\\
 A_{q,j+1,r}^{(\eps)}
 &=-\frac1\al\sum_{n=0}^{r-1}
 nA_{q,j,n}^{(\eps)}
 \left(-\frac{\be_\eps}{\al}\right)^{r-n-1}.
 \label{eq:A-update}
\end{align}
Consequently,
\begin{equation}
 \boxed{\displaystyle
 d_r^{(\eps)}=
 \sum_{q=1}^{\lfloor(r+1)/2\rfloor}
 \frac{(-1)^q}{q!}A_{q,q-1,r}^{(\eps)}.}
 \label{eq:d-array}
\end{equation}
\Cref{eq:hqs,eq:A-initial,eq:A-update,eq:d-array} are a finite, nonrecursive-in-\(D\) coefficient engine.  They are particularly useful for exact rational arithmetic because every \(\lambda_r^{(\eps)}\) is rational and the only transcendental symbol is \(\al=\log4\).

\subsection{Generic low-order coefficients}

Suppressing the parity superscript and writing \(\be=\be_\eps\), \(\lambda_j=\lambda_j^{(\eps)}\), the first coefficients generated by any of the three equivalent mechanisms are
\begin{align}
 d_1={}&-\frac{\lambda_1}{\al},
 \label{eq:d1-general}\\
 d_2={}&-\frac{\lambda_2}{\al}+\frac{\be\lambda_1}{\al^2},
 \label{eq:d2-general}\\
 d_3={}&-\frac{\lambda_3}{\al}
 +\frac{\be\lambda_2-\lambda_1^2}{\al^2}
 -\frac{\be^2\lambda_1}{\al^3},
 \label{eq:d3-general}\\
 d_4={}&-\frac{\lambda_4}{\al}
 +\frac{\be\lambda_3-3\lambda_1\lambda_2}{\al^2}
 +\frac{-\be^2\lambda_2+\tfrac52\be\lambda_1^2}{\al^3}
 +\frac{\be^3\lambda_1}{\al^4},
 \label{eq:d4-general}\\
 d_5={}&-\frac{\lambda_5}{\al}
 +\frac{\be\lambda_4-4\lambda_1\lambda_3-2\lambda_2^2}{\al^2}
 +\frac{-\be^2\lambda_3+7\be\lambda_1\lambda_2-2\lambda_1^3}{\al^3}
 \nonumber\\
 &\quad+\frac{\be^3\lambda_2-\tfrac92\be^2\lambda_1^2}{\al^4}
 -\frac{\be^4\lambda_1}{\al^5}.
 \label{eq:d5-general}
\end{align}
These formulae expose the combinatorial structure: products of logarithmic coefficients record set partitions generated by composition, while powers of \(\be\) come from repeated differentiation of the logarithmic carrier.

\subsection{Specialized coefficients for the two swing sheets}

Substituting \(\be_0=-1/2\) and the even \(\lambda_r\)'s gives
\begin{align}
 d_1^{(0)}&=\frac{1}{8\al},
 &d_2^{(0)}&=\frac{1}{16\al^2},
 \label{eq:d-even-12}\\
 d_3^{(0)}&=-\frac{\al^2+3\al-6}{192\al^3},
 &d_4^{(0)}&=-\frac{2\al^2+15\al-12}{768\al^4},
 \label{eq:d-even-34}\\
 d_5^{(0)}&=\frac{12\al^4+20\al^3+20\al^2-135\al+60}{7680\al^5},
 \label{eq:d-even-5}\\
 d_6^{(0)}&=\frac{24\al^4+90\al^3+215\al^2-420\al+120}{30720\al^6}.
 \label{eq:d-even-6}
\end{align}
For the odd sheet,
\begin{align}
 d_1^{(1)}&=-\frac{3}{8\al},
 &d_2^{(1)}&=\frac{2\al+3}{16\al^2},
 \label{eq:d-odd-12}\\
 d_3^{(1)}&=-\frac{3\al^2+13\al+6}{64\al^3},
 &d_4^{(1)}&=\frac{4\al^3+42\al^2+53\al+12}{256\al^4},
 \label{eq:d-odd-34}\\
 d_5^{(1)}&=-\frac{12\al^4+280\al^3+720\al^2+445\al+60}{2560\al^5},
 \label{eq:d-odd-5}\\
 d_6^{(1)}&=\frac{80\al^5+1872\al^4+9030\al^3+10845\al^2+4020\al+360}{30720\al^6}.
 \label{eq:d-odd-6}
\end{align}

\begin{theorem}[All-orders inverse transseries]
\label{thm:inverse-main}
For every fixed \(R\ge0\), as \(y\to+\infty\),
\begin{equation}
 \boxed{\displaystyle
 M_\eps(y)=w_\eps(y)+\sum_{r=1}^R
 \frac{d_r^{(\eps)}}{w_\eps(y)^r}
 +\OO\!\left(w_\eps(y)^{-R-1}\right),}
 \label{eq:inverse-fixed-R}
\end{equation}
where \(w_0,w_1\) are given by \cref{eq:w0,eq:w1}.  Equivalently, for the original index,
\begin{equation}
 \boxed{\displaystyle
 N_\eps(y)=2w_\eps(y)+\eps+2\sum_{r=1}^R
 \frac{d_r^{(\eps)}}{w_\eps(y)^r}
 +\OO\!\left(w_\eps(y)^{-R-1}\right).}
 \label{eq:N-inverse-fixed-R}
\end{equation}
\end{theorem}

\begin{proof}
The formal construction cancels the residual in \cref{eq:full-log-inverse-equation} through order \(w^{-R}\), leaving \(\OO(w^{-R-1})\).  The derivative of the full logarithmic phase is
\begin{equation*}
 \al+\frac{\be_\eps}{m}+\OO(m^{-2})=\al+o(1),
\end{equation*}
so it is bounded away from zero for large \(m\).  The mean-value theorem therefore converts the phase residual into an inverse-variable error of the same order.  Finally \(N_\eps=2M_\eps+\eps\).
\end{proof}

\begin{remark}[Residual certification]
A robust implementation should not trust a long symbolic expression merely because it looks patterned.  Given a proposed truncation \(D_R(t)\), substitute it into the left-hand side of \cref{eq:fundamental-D}; the coefficients of \(t,t^2,\ldots,t^R\) must vanish identically.  This supplies an exact algebraic certificate for every table of \(d_r\)'s.
\end{remark}

\section{Unfolding the Lambert core: polynomial--logarithmic transseries}
\label{sec:polylog}

The Lambert-normalized expansion is usually the most accurate form for computation because it resums the entire dominant power--logarithmic phase.  For comparison with classical inverse asymptotics, however, it is useful to expand everything in powers of \(1/\log y\) with polynomial coefficients in \(\log\log y\).

\subsection{Normalized variables}

Set
\begin{equation}
 x=\al m,
 \qquad
 X_\eps=L_\eps(y)+\be_\eps\log\al
 =\log\!\left(\frac{y\al^{\be_\eps}}{K_\eps}\right),
 \qquad
 \ell_\eps=\log X_\eps.
 \label{eq:X-def}
\end{equation}
Explicitly,
\begin{equation}
 X_0=\log\!\left(y\sqrt{\frac{\pi}{\al}}\right),
 \qquad
 X_1=\log\!\left(\frac{y\sqrt{\pi\al}}{2}\right).
 \label{eq:X-explicit}
\end{equation}
Define scaled logarithmic coefficients
\begin{equation}
 \mu_k^{(\eps)}=\al^k\lambda_k^{(\eps)}.
 \label{eq:mu-def}
\end{equation}
Then \cref{eq:full-log-inverse-equation} becomes
\begin{equation}
 X_\eps=x+\be_\eps\log x+
 \sum_{k\ge1}\mu_k^{(\eps)}x^{-k}.
 \label{eq:x-normalized-equation}
\end{equation}

We seek
\begin{equation}
 x\sim X-\be\ell+\sum_{r\ge1}\frac{P_r(\ell)}{X^r},
 \qquad \ell=\log X,
 \label{eq:P-ansatz}
\end{equation}
where the parity labels are temporarily suppressed.

\subsection{All-orders polynomial recurrence}

Let \(u=X^{-1}\) and, for \(r\ge1\), define
\begin{equation}
 R_{<r}(u,\ell)=-\be\ell+
 \sum_{j=1}^{r-1}P_j(\ell)u^j.
 \label{eq:R-polylog}
\end{equation}
Substituting \(x=X+R\) into \cref{eq:x-normalized-equation} gives
\begin{equation}
 R+\be\ell+
 \be\log(1+uR)+
 \sum_{k\ge1}\mu_k\left(\frac{u}{1+uR}\right)^k=0.
 \label{eq:R-polylog-equation}
\end{equation}

\begin{theorem}[Polynomial--logarithmic coefficient recurrence]
\label{thm:P-recurrence}
The polynomials \(P_r(\ell)\) are uniquely determined by
\begin{equation}
 \boxed{\displaystyle
 P_r(\ell)=-\Coeff{u^r}{
 \be\log\!\left(1+uR_{<r}(u,\ell)\right)
 +\sum_{k=1}^r\mu_k
 \left(\frac{u}{1+uR_{<r}(u,\ell)}\right)^k}.}
 \label{eq:P-recurrence}
\end{equation}
For every \(r\ge1\), \(P_r\) has degree at most \(r\), and because \(\be\ne0\) in both swing-factorial sheets,
\begin{equation}
 \Coeff{\ell^r}{P_r(\ell)}=\frac{\be^{r+1}}{r}.
 \label{eq:P-leading}
\end{equation}
Thus every scalar coefficient \(p_{r,j}=\Coeff{\ell^j}{P_r(\ell)}\) has the finite general formula obtained by applying \(\Coeff{u^r\ell^j}{}\) to the right-hand side of \cref{eq:P-recurrence}.
\end{theorem}

\begin{proof}
In the coefficient of \(u^r\) in the nonlinear terms of \cref{eq:R-polylog-equation}, \(P_r\) cannot occur because it would enter through \(uR\) and hence at order at least \(u^{r+1}\).  The \(u^r\) coefficient of \(R+\be\ell\) is exactly \(P_r\), proving the recurrence and uniqueness.

Assume inductively that \(\deg P_j\le j\) for \(j<r\).  Every occurrence of \(R_{<r}\) contributes at most one power of \(\ell\) per power of \(u\), so the coefficient of \(u^r\) has degree at most \(r\).  The terms involving \(\mu_k\) lose at least \(k\ge1\) powers of the available \(uR\) factors and therefore have degree at most \(r-1\).  The top-degree coefficient is consequently the same as in the pure equation \(X=x+\be\log x\), whose exact Lambert solution, or a direct induction in \cref{eq:P-recurrence}, gives \(\be^{r+1}/r\).
\end{proof}

\subsection{Generic first polynomials}

The recurrence gives
\begin{align}
 P_1(\ell)={}&\be^2\ell-\mu_1,
 \label{eq:P1-general}\\
 P_2(\ell)={}&\frac{\be^3}{2}\ell^2
 -(\be^3+\be\mu_1)\ell
 +\be\mu_1-\mu_2,
 \label{eq:P2-general}\\
 P_3(\ell)={}&\frac{\be^4}{3}\ell^3
 -\left(\frac32\be^4+\be^2\mu_1\right)\ell^2
 +\left(\be^4+3\be^2\mu_1-2\be\mu_2\right)\ell
 \nonumber\\
 &\quad-\be^2\mu_1+\be\mu_2-\mu_1^2-\mu_3,
 \label{eq:P3-general}\\
 P_4(\ell)={}&\frac{\be^5}{4}\ell^4
 -\left(\frac{11}{6}\be^5+\be^3\mu_1\right)\ell^3
 +\left(3\be^5+\frac{11}{2}\be^3\mu_1-3\be^2\mu_2\right)\ell^2
 \nonumber\\
 &\quad+\left(-\be^5-6\be^3\mu_1+5\be^2\mu_2
 -3\be\mu_1^2-3\be\mu_3\right)\ell
 \nonumber\\
 &\quad+\be^3\mu_1-\be^2\mu_2+\frac52\be\mu_1^2
 +\be\mu_3-3\mu_1\mu_2-\mu_4.
 \label{eq:P4-general}
\end{align}
These expressions separate the universal Lambert contribution (the pure powers of \(\be\)) from the gamma-ratio corrections (the terms containing at least one \(\mu_k\)).

\subsection{Specialized logarithmic expansions}

For the even sheet, the first three correction polynomials are
\begin{align}
 P_1^{(0)}(\ell)&=\frac14\ell+\frac\al8,
 \label{eq:P-even-1}\\
 P_2^{(0)}(\ell)&=-\frac1{16}\ell^2
 +\left(\frac18-\frac\al{16}\right)\ell+\frac\al{16},
 \label{eq:P-even-2}\\
 P_3^{(0)}(\ell)&=\frac1{48}\ell^3
 +\left(\frac\al{32}-\frac3{32}\right)\ell^2
 +\left(\frac1{16}-\frac{3\al}{32}\right)\ell
 -\frac{\al^3}{192}-\frac{\al^2}{64}+\frac\al{32}.
 \label{eq:P-even-3}
\end{align}
Thus
\begin{equation}
 M_0(y)\sim\frac1\al\left[
 X_0+\frac12\ell_0+
 \frac{P_1^{(0)}(\ell_0)}{X_0}+
 \frac{P_2^{(0)}(\ell_0)}{X_0^2}+
 \frac{P_3^{(0)}(\ell_0)}{X_0^3}+\cdots\right].
 \label{eq:M0-polylog}
\end{equation}
For the odd sheet,
\begin{align}
 P_1^{(1)}(\ell)&=\frac14\ell-\frac{3\al}{8},
 \label{eq:P-odd-1}\\
 P_2^{(1)}(\ell)&=\frac1{16}\ell^2
 -\left(\frac18+\frac{3\al}{16}\right)\ell
 +\frac{\al^2}{8}+\frac{3\al}{16},
 \label{eq:P-odd-2}\\
 P_3^{(1)}(\ell)&=\frac1{48}\ell^3
 -\left(\frac3{32}+\frac{3\al}{32}\right)\ell^2
 +\left(\frac1{16}+\frac{9\al}{32}+\frac{\al^2}{8}\right)\ell
 \nonumber\\
 &\quad-\frac{3\al}{32}-\frac{13\al^2}{64}-\frac{3\al^3}{64},
 \label{eq:P-odd-3}
\end{align}
and
\begin{equation}
 M_1(y)\sim\frac1\al\left[
 X_1-\frac12\ell_1+
 \frac{P_1^{(1)}(\ell_1)}{X_1}+
 \frac{P_2^{(1)}(\ell_1)}{X_1^2}+
 \frac{P_3^{(1)}(\ell_1)}{X_1^3}+\cdots\right].
 \label{eq:M1-polylog}
\end{equation}

\begin{corollary}[Leading elementary forms]
As \(y\to\infty\),
\begin{align}
 M_0(y)&=\frac1\al\left(
 X_0+\frac12\log X_0+\OO\!\left(\frac{\log X_0}{X_0}\right)\right),
 \label{eq:M0-leading}\\
 M_1(y)&=\frac1\al\left(
 X_1-\frac12\log X_1+\OO\!\left(\frac{\log X_1}{X_1}\right)\right).
 \label{eq:M1-leading}
\end{align}
\end{corollary}

The opposite signs of \(\tfrac12\log X\) are the inverse-side signature of the direct powers \(m^{-1/2}\) and \(m^{1/2}\).

\subsection{Equivalence of the Lambert and logarithmic normalizations}

The two inverse expansions do not contain different information.  The Lambert form first solves the carrier exactly and then expands only the small gamma-ratio correction.  The polynomial--logarithmic form expands the Lambert carrier itself.  More concretely, expanding \(w_\eps(y)\) in the scale \(X^{-1}\R[\ell]\) and then substituting that expansion into \cref{eq:inverse-fixed-R} reproduces \cref{eq:P-recurrence} coefficient by coefficient.  Conversely, resumming the subfamily of terms that survives when all \(\mu_k\) vanish reconstructs the appropriate Lambert branch.

\begin{summarybox}[Which inverse expansion should be used?]
For numerical work, use the Lambert-normalized series \cref{eq:inverse-fixed-R}; it absorbs the largest logarithmic effects and usually needs fewer terms.  For symbolic comparison, scale classification, or interaction with a polynomial--logarithmic transseries algebra, use \cref{eq:M0-polylog,eq:M1-polylog} and recurrence \cref{eq:P-recurrence}.
\end{summarybox}

\section{Beyond all algebraic orders}
\label{sec:beyond}

The coefficient formulae above completely determine the Poincar\'e transseries, but the Bernoulli late terms show that this series cannot be summed by taking infinitely many terms at a fixed \(m\).  There are two clean ways to retain the missing information.

\subsection{Exact completion by Binet's integral}

\Cref{eq:exact-Binet-swing} is already an exact completion.  For a chosen truncation order \(N\), define the exact logarithmic remainder
\begin{equation}
 \mathcal R_{\eps,N}(m)=
 \log S_\eps(m)-\log K_\eps-\al m-\be_\eps\log m
 -\sum_{r=1}^{N-1}\lambda_r^{(\eps)}m^{-r}.
 \label{eq:exact-remainder}
\end{equation}
It can be evaluated either from gamma functions or from the difference of the two Binet integrals.  For fixed \(N\), \(\mathcal R_{\eps,N}(m)=\OO(m^{-N})\).  If \(N\) is selected near the least term, the remainder is exponentially small in the standard gamma-function sense.  This exact remainder is preferable to an unqualified symbol such as ``\(+Ce^{-2\pi m}\)'', because the most useful exponentially improved representations involve terminant functions and Stokes smoothing rather than a single universal constant \(C\).

\subsection{Transport of a general exponentially small sector}

The inverse transport mechanism is independent of how the beyond-all-orders remainder is represented.

\begin{theorem}[Transseries-sector transport through inversion]
\label{thm:sector-transport}
Suppose an exact or exponentially improved forward phase is written
\begin{equation}
 Y=\Phi_{\mathrm{alg}}(m)+E(m),
 \label{eq:sector-forward}
\end{equation}
where \(\Phi_{\mathrm{alg}}'(m)\ne0\) and \(E\) is smaller than the algebraic scale retained in \(\Phi_{\mathrm{alg}}\).  Let \(m_0\) solve \(\Phi_{\mathrm{alg}}(m_0)=Y\).  Formally in the small sector \(E\),
\begin{equation}
 m=m_0+\sum_{q\ge1}\frac{(-1)^q}{q!}
 \left(\frac1{\Phi_{\mathrm{alg}}'(m_0)}\frac{d}{dm_0}\right)^{q-1}
 \left(\frac{E(m_0)^q}{\Phi_{\mathrm{alg}}'(m_0)}\right).
 \label{eq:sector-transport-all}
\end{equation}
In particular,
\begin{align}
 m={}&m_0-\frac{E}{\Phi'}
 +\frac{EE'}{(\Phi')^2}
 -\frac{\Phi''E^2}{2(\Phi')^3}
 +\OO(E^3,E^2E',E(E')^2),
 \label{eq:sector-transport-first}
\end{align}
where the functions on the right are evaluated at \(m_0\).
\end{theorem}

\begin{proof}
This is \cref{lem:perturbative-LB} with \(F_0=\Phi_{\mathrm{alg}}\), \(h=E\), and the coupling parameter set to one after formal expansion.  The displayed first terms follow by applying the differential operator once.
\end{proof}

For the swing factorial, \(\Phi_{\mathrm{alg}}'(m)=\al+\OO(m^{-1})\).  Therefore the leading transport law is especially simple:
\begin{equation}
 \Delta M_\eps(y)\sim-\frac{E(w_\eps(y))}{\al}.
 \label{eq:swing-sector-leading}
\end{equation}
If a chosen exponentially improved gamma representation has a sector of magnitude \(m^\rho e^{-\kappa m}\), its inverse image has magnitude \(w_\eps(y)^\rho e^{-\kappa w_\eps(y)}/\al\) at first order, while products in \cref{eq:sector-transport-all} generate the higher multiples of that sector.  Since \(w_\eps(y)\asymp(\log y)/\al\), an exponential in \(w\) becomes an additional power-like scale in \(y\); for example,
\begin{equation}
 e^{-2\pi w_\eps(y)}
 =y^{-2\pi/\al}\,(\log y)^{\OO(1)}
 \times\hbox{an explicitly expandable constant/logarithmic factor}.
 \label{eq:exp-to-y-power}
\end{equation}
This observation explains why hyperasymptotic corrections to the inverse form a genuinely finer transseries grid than the algebraic powers of \(1/\log y\).

\section{Discrete inverses and index recovery}
\label{sec:discrete}

The analytic inverse \(M_\eps(y)\) answers the smooth reversion problem.  For an integer-valued inverse one usually wants the first term on a fixed parity sheet that reaches a prescribed level.

\begin{definition}[Branchwise upper inverse]
For \(y\ge1\), define
\begin{equation}
 J_\eps(y)=\min\{m\in\N_0:S_\eps(m)\ge y\},
 \qquad
 \nu_\eps(y)=2J_\eps(y)+\eps.
 \label{eq:discrete-inverse}
\end{equation}
\end{definition}

By monotonicity,
\begin{equation}
 J_\eps(y)=\left\lceil M_\eps(y)\right\rceil,
 \qquad
 \nu_\eps(y)=2\left\lceil M_\eps(y)\right\rceil+\eps.
 \label{eq:ceil-inverse}
\end{equation}
This identity includes the case in which \(y\) is exactly a branch value, because the ceiling of the corresponding integral inverse is unchanged.

In floating-point computation, one should use the asymptotic expansion only to produce a candidate integer.  The final answer is certified by comparing \(\log S_\eps(m)\) with \(\log y\) at the candidate and its neighbor.  This avoids a rounding error when \(y\) lies exponentially close to an exact branch value.

The first index of either parity that reaches \(y\) is
\begin{equation}
 \nu_{\min}(y)=\min\{\nu_0(y),\nu_1(y)\}.
 \label{eq:global-first-hit}
\end{equation}
Although this is a well-defined generalized inverse, it does not make the original sequence monotone after the first hit.  The odd sheet eventually wins by an increasing margin.  Indeed, from \cref{eq:M0-leading,eq:M1-leading},
\begin{equation}
 N_0(y)-N_1(y)
 =\frac{2}{\al}\log\log y
 +\frac{2}{\al}\log\frac{2}{\al}-1+o(1),
 \qquad y\to\infty.
 \label{eq:sheet-gap}
\end{equation}
Thus large levels are first reached on the odd sheet, roughly \((2/\log4)\log\log y\) indices before the even sheet catches up.

\section{Exact-arithmetic generation and residual tests}
\label{sec:algorithm}

All coefficients can be generated without numerical approximation.  The following compact SymPy-style code mirrors the mathematical recurrences; \texttt{alpha} is kept symbolic and Bernoulli values are exact rationals.

\begin{lstlisting}[style=code,language=Python,caption={Exact generation of direct and inverse coefficients.}]
from sympy import (Symbol, Rational, bernoulli, log,
                   series, expand, simplify)

t = Symbol('t')
alpha = Symbol('alpha', nonzero=True)

def lam(eps, k):
    a = Rational(eps) + Rational(1, 2)
    return simplify(
        (-1)**(k + 1) * (bernoulli(k + 1, a) - bernoulli(k + 1, 1))
        / (k * (k + 1))
    )

def direct_coefficients(eps, order):
    c = [Rational(1)]
    for r in range(1, order + 1):
        c.append(simplify(sum(k * lam(eps, k) * c[r-k]
                              for k in range(1, r+1)) / r))
    return c

def inverse_coefficients(eps, order):
    beta = Rational(eps) - Rational(1, 2)
    H = sum(lam(eps, k) * t**k for k in range(1, order + 1))
    D = 0
    d = []
    for r in range(1, order + 1):
        u = t / (1 + t*D)
        nonlinear = beta*log(1 + t*D) + H.subs(t, u)
        dr = -series(nonlinear, t, 0, r+1).removeO().expand().coeff(t, r) / alpha
        dr = simplify(dr)
        d.append(dr)
        D = expand(D + dr*t**r)

        # Algebraic certificate: all terms through t^r vanish.
        residual = alpha*D + beta*log(1+t*D) + H.subs(t, t/(1+t*D))
        assert series(residual, t, 0, r+1).removeO().expand() == 0
    return d
\end{lstlisting}

The computational workflow is therefore:
\begin{enumerate}
\item generate \(\lambda_r^{(\eps)}\) from \cref{eq:lambda-general};
\item generate \(c_r^{(\eps)}\) from \cref{eq:c-recurrence};
\item generate \(d_r^{(\eps)}\) from \cref{eq:d-recurrence} or \cref{eq:d-array};
\item verify the formal residual in \cref{eq:fundamental-D};
\item for a numerical \(y\), evaluate the appropriate Lambert core and add the desired number of inverse corrections;
\item if a discrete answer is needed, certify neighboring integers using log-gamma evaluations.
\end{enumerate}

\section{Numerical validation}
\label{sec:numerics}

The following tables use high-precision gamma and Lambert evaluations.  They are not proofs; their role is to expose normalization, sign, or branch errors that formal manipulations might otherwise conceal.

\subsection{Direct expansion}

Let
\begin{equation}
 S_{\eps,R}^{\mathrm{dir}}(m)=K_\eps4^m m^{\be_\eps}
 \sum_{r=0}^{R}c_r^{(\eps)}m^{-r}.
 \label{eq:direct-trunc}
\end{equation}
\Cref{tab:direct-errors} reports the relative error
\(\left|S_{\eps,R}^{\mathrm{dir}}/S_\eps-1\right|\).

\begin{table}[htbp]
\centering
\caption{Relative errors of the direct expansion.  Column \(R\) includes terms through \(c_Rm^{-R}\).}
\label{tab:direct-errors}
\small
\begin{tabular}{@{}ccrrrr@{}}
\toprule
\(\eps\)&\(m\)&\(R=0\)&\(R=2\)&\(R=4\)&\(R=8\)\\
\midrule
0&10  &\(1.2573\!\times10^{-2}\)&\(4.8642\!\times10^{-6}\)&\(1.5087\!\times10^{-8}\)&\(1.6336\!\times10^{-12}\)\\
0&50  &\(2.5031\!\times10^{-3}\)&\(3.9053\!\times10^{-8}\)&\(4.8680\!\times10^{-12}\)&\(8.5498\!\times10^{-19}\)\\
0&200 &\(6.2519\!\times10^{-4}\)&\(6.1033\!\times10^{-10}\)&\(4.7561\!\times10^{-15}\)&\(3.2645\!\times10^{-24}\)\\
\addlinespace
1&10  &\(3.5645\!\times10^{-2}\)&\(8.6312\!\times10^{-6}\)&\(1.8177\!\times10^{-8}\)&\(1.7087\!\times10^{-12}\)\\
1&50  &\(7.4227\!\times10^{-3}\)&\(7.0071\!\times10^{-8}\)&\(5.8858\!\times10^{-12}\)&\(8.9455\!\times10^{-19}\)\\
1&200 &\(1.8701\!\times10^{-3}\)&\(1.0977\!\times10^{-9}\)&\(5.7556\!\times10^{-15}\)&\(3.4163\!\times10^{-24}\)\\
\bottomrule
\end{tabular}
\end{table}

\subsection{Inverse expansion}

Choose \(y=S_\eps(m)\), so the exact inverse is the displayed \(m\), and define
\begin{equation}
 M_{\eps,R}^{\mathrm{inv}}(y)=w_\eps(y)+
 \sum_{r=1}^{R}d_r^{(\eps)}w_\eps(y)^{-r}.
 \label{eq:inverse-trunc}
\end{equation}
\Cref{tab:inverse-errors} gives the absolute error
\(\left|M_{\eps,R}^{\mathrm{inv}}(S_\eps(m))-m\right|\).

\begin{table}[htbp]
\centering
\caption{Absolute errors of the Lambert-normalized inverse expansion.  The ``core'' column is \(R=0\).}
\label{tab:inverse-errors}
\small
\begin{tabular}{@{}ccrrrr@{}}
\toprule
\(\eps\)&\(m\)&core&\(R=2\)&\(R=4\)&\(R=8\)\\
\midrule
0&10  &\(9.3505\!\times10^{-3}\)&\(6.0195\!\times10^{-7}\)&\(3.4687\!\times10^{-9}\)&\(1.6428\!\times10^{-12}\)\\
0&50  &\(1.8164\!\times10^{-3}\)&\(1.9743\!\times10^{-9}\)&\(8.0457\!\times10^{-13}\)&\(8.0953\!\times10^{-19}\)\\
0&200 &\(4.5166\!\times10^{-4}\)&\(2.2503\!\times10^{-11}\)&\(7.2653\!\times10^{-16}\)&\(3.0590\!\times10^{-24}\)\\
\addlinespace
1&10  &\(2.5271\!\times10^{-2}\)&\(1.5678\!\times10^{-4}\)&\(1.9145\!\times10^{-6}\)&\(5.0787\!\times10^{-10}\)\\
1&50  &\(5.3358\!\times10^{-3}\)&\(1.3679\!\times10^{-6}\)&\(6.7900\!\times10^{-10}\)&\(2.9371\!\times10^{-16}\)\\
1&200 &\(1.3478\!\times10^{-3}\)&\(2.1720\!\times10^{-8}\)&\(6.7553\!\times10^{-13}\)&\(1.1441\!\times10^{-21}\)\\
\bottomrule
\end{tabular}
\end{table}

The odd-sheet inverse converges more slowly at the same \(m\) in these low-order tests because its logarithmic correction begins with \(+3/(8m)\), three times the magnitude of the even sheet's first correction.  Nevertheless the all-orders pattern and the predicted powers of \(w^{-1}\) are clearly visible.

\begin{warningbox}[Optimal truncation]
The decreasing errors in these tables do not imply convergence as \(R\to\infty\) at fixed \(m\).  The direct Bernoulli series and the induced inverse series are asymptotic.  Terms should be stopped near their least magnitude; Binet's exact integral or a terminant expansion should be used beyond that point.
\end{warningbox}

\section{A symmetric centered normal form}
\label{sec:centered}

The uncentered coordinate \(m\) is natural from the sequence definition, but the gamma ratio possesses a better symmetric coordinate in which every odd inverse power disappears from the logarithmic correction.

\subsection{Common centering of both parity sheets}

Define
\begin{equation}
 s_\eps=\frac{2\eps+1}{4},
 \qquad x=m+s_\eps.
 \label{eq:center-shift}
\end{equation}
Then \(s_0=1/4\), \(s_1=3/4\), and
\begin{align}
 m+\eps+\frac12&=x+s_\eps,
 &m+1&=x+1-s_\eps,
 &\be_\eps&=2s_\eps-1.
 \label{eq:center-identities}
\end{align}
Moreover,
\begin{equation}
 K_\eps4^{-s_\eps}=\frac1{\sqrt{2\pi}}
 \label{eq:center-common-K}
\end{equation}
for both values of \(\eps\).  Hence the exact common centered representation is
\begin{equation}
 \boxed{\displaystyle
 S_\eps(m)=\frac{e^{\al x}}{\sqrt{2\pi}}
 \frac{\Gamma(x+s_\eps)}{\Gamma(x+1-s_\eps)},
 \qquad x=m+s_\eps.}
 \label{eq:center-exact}
\end{equation}
This simultaneously equalizes the leading constant and makes the two gamma shifts complementary.

\subsection{Even-power logarithmic series}

Applying shifted Stirling in the coordinate \(x\) gives
\begin{equation}
 \log S_\eps(m)
 \sim -\frac12\log(2\pi)+\al x+\be_\eps\log x
 +\sum_{r\ge1}\widehat\lambda_r^{(\eps)}x^{-r},
 \label{eq:center-log}
\end{equation}
where
\begin{equation}
 \widehat\lambda_r^{(\eps)}=
 \frac{(-1)^{r+1}}{r(r+1)}
 \left[\B_{r+1}(s_\eps)-\B_{r+1}(1-s_\eps)\right].
 \label{eq:center-lambda-general}
\end{equation}
The Bernoulli reflection identity
\begin{equation}
 \B_n(1-z)=(-1)^n\B_n(z)
 \label{eq:Bernoulli-reflection}
\end{equation}
therefore implies
\begin{align}
 \widehat\lambda_{2q-1}^{(\eps)}&=0,
 \label{eq:center-odd-zero}\\
 \widehat\lambda_{2q}^{(\eps)}
 &=-\frac{\B_{2q+1}(s_\eps)}{q(2q+1)},
 \qquad q\ge1.
 \label{eq:center-even-coeff}
\end{align}
Since \(s_1=1-s_0\),
\begin{equation}
 \widehat\lambda_{2q}^{(1)}=-\widehat\lambda_{2q}^{(0)}.
 \label{eq:center-lambda-sign}
\end{equation}
The first nonzero values are
\begin{align}
 \widehat\lambda_2^{(0)}&=-\frac1{64},&
 \widehat\lambda_4^{(0)}&=\frac5{2048},&
 \widehat\lambda_6^{(0)}&=-\frac{61}{49152},&
 \widehat\lambda_8^{(0)}&=\frac{1385}{1048576},
 \label{eq:center-lambda-even-first}
\end{align}
with opposite signs for \(\eps=1\).

Let
\begin{equation}
 \widehat H_\eps(z)=\sum_{q\ge1}\widehat\lambda_{2q}^{(\eps)}z^q,
 \qquad
 e^{\widehat H_\eps(z)}=\sum_{q\ge0}\widehat c_q^{(\eps)}z^q.
 \label{eq:center-H-c}
\end{equation}
Then
\begin{equation}
 \widehat c_q^{(\eps)}=
 \frac1{q!}Y_q\!\left(
 1!\widehat\lambda_2^{(\eps)},
 2!\widehat\lambda_4^{(\eps)},\ldots,
 q!\widehat\lambda_{2q}^{(\eps)}
 \right)
 \label{eq:center-c-Bell}
\end{equation}
and
\begin{equation}
 \boxed{\displaystyle
 S_\eps(m)\sim\frac{e^{\al x}x^{\be_\eps}}{\sqrt{2\pi}}
 \sum_{q\ge0}\widehat c_q^{(\eps)}x^{-2q}.}
 \label{eq:center-direct}
\end{equation}
The first two centered expansions are
\begin{align}
 S_0(m)&\sim\frac{e^{\al x}}{\sqrt{2\pi x}}
 \left(1-\frac1{64x^2}+\frac{21}{8192x^4}
 -\frac{671}{524288x^6}+\frac{180323}{134217728x^8}+\cdots\right),
 \quad x=m+\frac14,
 \label{eq:center-even-direct}\\
 S_1(m)&\sim\frac{e^{\al x}\sqrt{x}}{\sqrt{2\pi}}
 \left(1+\frac1{64x^2}-\frac{19}{8192x^4}
 +\frac{631}{524288x^6}-\frac{174317}{134217728x^8}+\cdots\right),
 \quad x=m+\frac34.
 \label{eq:center-odd-direct}
\end{align}
The centered form is not a new asymptotic function; it is a resummation of the same formal series after a translation.  Its advantage is structural sparsity.

\subsection{Centered Lambert carriers and inverse expansion}

The centered carrier equation is
\begin{equation}
 y=\frac{e^{\al\widehat w}\widehat w^{\be_\eps}}{\sqrt{2\pi}}.
 \label{eq:center-carrier-equation}
\end{equation}
Its large real solutions are
\begin{align}
 \widehat w_0(y)&=-\frac1{2\al}\W_{-1}\!\left(-\frac{\al}{\pi y^2}\right),
 \label{eq:center-w0}\\
 \widehat w_1(y)&=\frac1{2\al}\W_0\!\left(4\pi\al y^2\right).
 \label{eq:center-w1}
\end{align}
Using the same operator calculus as in \cref{sec:inverse-coefficients}, but replacing \(H_\eps(t)\) by
\begin{equation}
 \widehat H_\eps^{\mathrm{inv}}(t)
 =\sum_{q\ge1}\widehat\lambda_{2q}^{(\eps)}t^{2q},
 \label{eq:center-H-inv}
\end{equation}
produces
\begin{equation}
 \boxed{\displaystyle
 M_\eps(y)\sim\widehat w_\eps(y)-s_\eps
 +\sum_{r\ge1}\widehat d_r^{(\eps)}\widehat w_\eps(y)^{-r}.}
 \label{eq:center-inverse}
\end{equation}
Every formula \cref{eq:d-recurrence,eq:d-closed,eq:d-array} remains valid after placing hats on \(H,d,w\).  Because \(\widehat H\) starts at \(t^2\),
\begin{equation}
 \widehat d_1^{(\eps)}=0,
 \label{eq:center-d1-zero}
\end{equation}
and the first coefficients are
\begin{align}
 \widehat d_2^{(0)}&=\frac1{64\al},&
 \widehat d_3^{(0)}&=\frac1{128\al^2},&
 \widehat d_4^{(0)}&=-\frac{5\al^2-8}{2048\al^3},
 \label{eq:center-d-even-first}\\
 \widehat d_2^{(1)}&=-\frac1{64\al},&
 \widehat d_3^{(1)}&=\frac1{128\al^2},&
 \widehat d_4^{(1)}&=\frac{5\al^2-8}{2048\al^3}.
 \label{eq:center-d-odd-first}
\end{align}
At the next orders,
\begin{align}
 \widehat d_5^{(0)}&=-\frac{7\al^2-8}{4096\al^4},&
 \widehat d_6^{(0)}&=\frac{61\al^4-57\al^2+48}{49152\al^5},
 \label{eq:center-d-even-next}\\
 \widehat d_5^{(1)}&=-\frac{7\al^2-8}{4096\al^4},&
 \widehat d_6^{(1)}&=-\frac{61\al^4-57\al^2+48}{49152\al^5}.
 \label{eq:center-d-odd-next}
\end{align}

\begin{table}[htbp]
\centering
\caption{Absolute inverse errors using only the centered carrier \(\widehat w_\eps-s_\eps\) and using centered corrections through \(R=4\).}
\label{tab:centered-inverse-errors}
\small
\begin{tabular}{@{}ccrr@{}}
\toprule
\(\eps\)&\(m\)&centered core&centered \(R=4\)\\
\midrule
0&10  &\(1.1103\!\times10^{-4}\)&\(2.6060\!\times10^{-9}\)\\
0&50  &\(4.4957\!\times10^{-6}\)&\(1.0842\!\times10^{-12}\)\\
0&200 &\(2.8158\!\times10^{-7}\)&\(1.1092\!\times10^{-15}\)\\
\addlinespace
1&10  &\(9.4239\!\times10^{-5}\)&\(2.9096\!\times10^{-9}\)\\
1&50  &\(4.3450\!\times10^{-6}\)&\(1.1084\!\times10^{-12}\)\\
1&200 &\(2.7917\!\times10^{-7}\)&\(1.1154\!\times10^{-15}\)\\
\bottomrule
\end{tabular}
\end{table}

The improvement over the uncentered core in \cref{tab:inverse-errors} is substantial: centering has absorbed the entire order-\(m^{-1}\) logarithmic correction into the carrier coordinate.  The centered and uncentered expansions are formally equivalent, but the centered one is often the preferable computational normal form.

\section{A reusable theorem for exponential gamma ratios}
\label{sec:general-gamma-ratio}

The swing factorial is an especially symmetric member of a wider class.  The derivation can be packaged once and reused.

\begin{theorem}[Gamma-ratio transseries and inverse]
\label{thm:general-gamma-ratio}
Let
\begin{equation}
 F(x)=K e^{\al x}\frac{\Gamma(x+a)}{\Gamma(x+b)},
 \qquad K>0,\quad \al>0,
 \label{eq:general-F}
\end{equation}
with fixed real \(a,b\) and \(\be=a-b\ne0\).  On any eventually increasing positive-real branch,
\begin{align}
 F(x)&\sim K e^{\al x}x^\be
 \exp\!\left(\sum_{r\ge1}\lambda_r x^{-r}\right),
 \label{eq:general-F-direct}\\
 \lambda_r&=\frac{(-1)^{r+1}}{r(r+1)}
 \left(\B_{r+1}(a)-\B_{r+1}(b)\right).
 \label{eq:general-F-lambda}
\end{align}
The multiplicative coefficients are given by \cref{eq:c-Bell,eq:c-partition,eq:c-recurrence}.  If
\begin{equation}
 w(y)=\frac\be\al\W_k\!\left(
 \frac\al\be\left(\frac yK\right)^{1/\be}
 \right)
 \label{eq:general-F-w}
\end{equation}
uses the Lambert branch for which \(w\to+\infty\), then
\begin{equation}
 F^{-1}(y)\sim w(y)+\sum_{r\ge1}d_rw(y)^{-r},
 \label{eq:general-F-inverse}
\end{equation}
where \(d_r\) is given by \cref{eq:d-recurrence,eq:d-closed,eq:d-array} with
\(
 H(t)=\sum_{r\ge1}\lambda_rt^r
\)
and \(P(t)=\al+\be t\).  The unfolded polynomials are given by \cref{eq:P-recurrence} with \(\mu_r=\al^r\lambda_r\).
\end{theorem}

\begin{proof}
The direct formula is the difference of two shifted Stirling expansions.  Bell exponentiation gives the multiplicative form.  The Lambert formula solves \(Ke^{\al w}w^\be=y\).  The inverse coefficient formulae follow from \cref{eq:fundamental-D,lem:perturbative-LB}; the unfolded recurrence follows from the normalization \(x\mapsto\al x\).  Eventual monotonicity guarantees that the selected real inverse branch is well-defined.
\end{proof}

\begin{remark}[Optimal centering in the general class]
Setting
\(
 s=(a+b-1)/2
\)
and \(X=x+s\) changes the shifts to \((1+\be)/2\) and \((1-\be)/2\), which sum to one.  Bernoulli reflection then eliminates every odd inverse power from the logarithmic gamma-ratio correction.  The swing-factorial shifts \(s_0=1/4\), \(s_1=3/4\) are precisely this general centering rule.
\end{remark}

\section{Further consequences}
\label{sec:consequences}

\subsection{Derivative and conditioning of the inverse}

Let
\begin{equation}
 \Phi_\eps(m)=\log S_\eps(m).
 \label{eq:Phi-def}
\end{equation}
The inverse-function theorem gives the exact identity
\begin{equation}
 M_\eps'(y)=\frac1{y\Phi_\eps'(M_\eps(y))}
 =\frac1{2y\left[\psi(2M_\eps+\eps+1)-\psi(M_\eps+1)\right]}.
 \label{eq:inverse-derivative-exact}
\end{equation}
Differentiating the logarithmic transseries yields
\begin{equation}
 \Phi_\eps'(m)\sim
 \al+\be_\eps m^{-1}-\sum_{r\ge1}r\lambda_r^{(\eps)}m^{-r-1}.
 \label{eq:Phi-prime-series}
\end{equation}
Define
\begin{equation}
 \frac1{\al+\be t-\sum_{r\ge1}r\lambda_rt^{r+1}}
 =\sum_{n\ge0}g_nt^n.
 \label{eq:g-def}
\end{equation}
Then
\begin{align}
 g_0&=\frac1\al,
 \label{eq:g0}\\
 g_n&=-\frac1\al\left(
 \be g_{n-1}-\sum_{r=1}^{n-1}r\lambda_rg_{n-r-1}
 \right),\qquad n\ge1.
 \label{eq:g-recurrence}
\end{align}
Consequently,
\begin{equation}
 M_\eps'(y)\sim\frac1y\sum_{n\ge0}
 g_n^{(\eps)}M_\eps(y)^{-n},
 \label{eq:Mprime-series}
\end{equation}
which becomes a Lambert-normalized expansion after substituting \cref{eq:inverse-fixed-R}.  In particular,
\begin{equation}
 M_\eps'(y)=\frac1{\al y}
 \left(1-\frac{\be_\eps}{\al M_\eps(y)}+
 \OO(M_\eps(y)^{-2})\right).
 \label{eq:Mprime-leading}
\end{equation}
Thus the absolute inverse sensitivity decays like \(1/(y\log4)\), while a relative perturbation in \(y\) produces an additive perturbation of order one in the branch parameter only after multiplication by \(1/\log4\).

\subsection{Complex inverse sheets}

For complex \(y\), after fixing branches of powers and logarithms, the carrier equation has the family
\begin{equation}
 w_{\eps,k}(y)=\frac{\be_\eps}{\al}
 \W_k\!\left(
 \frac{\al}{\be_\eps}
 \left(\frac{y}{K_\eps}\right)^{1/\be_\eps}
 \right),
 \qquad k\in\mathbb Z.
 \label{eq:complex-carriers}
\end{equation}
Whenever a selected branch satisfies \(|w_{\eps,k}|\to\infty\) in a sector avoiding gamma poles and logarithmic cuts, the same coefficients \(d_r^{(\eps)}\) give
\begin{equation}
 M_{\eps,k}(y)\sim w_{\eps,k}(y)+
 \sum_{r\ge1}d_r^{(\eps)}w_{\eps,k}(y)^{-r}.
 \label{eq:complex-inverse}
\end{equation}
The real inverses studied above are the distinguished choices \(k=-1\) for the even carrier approaching the negative cut from the real side and \(k=0\) for the positive odd carrier.

\subsection{Structural hierarchy of the coefficient families}

The complete calculation can be summarized as a chain of functorial coefficient transformations:
\begin{equation}
 \boxed{
 \text{Bernoulli polynomials}
 \xrightarrow{\text{gamma-ratio difference}}
 \lambda_r
 \xrightarrow{\exp/\text{Bell}}
 c_r
 \xrightarrow{\text{Lagrange--B\"urmann}}
 d_r
 \xrightarrow{\text{Lambert unfolding}}
 P_r(\ell).}
 \label{eq:coefficient-hierarchy}
\end{equation}
Every arrow is constructive at arbitrary order, and each has both a closed coefficient-extraction description and a triangular recurrence.  This is the precise sense in which the calculation belongs to a combinatorial coefficient calculus rather than to a finite list of hand-expanded Stirling terms.

\section{Conclusions}
\label{sec:conclusion}

The swing factorial has a richer asymptotic inverse structure than its elementary factorial quotient suggests.  The decisive facts are as follows.

\begin{enumerate}
\item The sequence is a period-two object.  Its even and odd subsequences are the gamma ratios \(S_0\) and \(S_1\), and the alternating drops rule out a single ordinary inverse.
\item The full direct logarithmic expansion has the Bernoulli-polynomial coefficient \cref{eq:lambda-general}.  Complete Bell polynomials, partition sums, and recurrence \cref{eq:c-recurrence} give every multiplicative coefficient.
\item The real inverse carriers are exact Lambert functions of different branches: \(\W_{-1}\) on the even sheet and \(\W_0\) on the odd sheet.
\item The inverse correction coefficients admit the triangular recurrence \cref{eq:d-recurrence}, the closed differential formula \cref{eq:d-closed}, and the finite array formula \cref{eq:d-array}.  These are general formulae at arbitrary order, not merely recursive numerical fits.
\item Unfolding the Lambert core gives the polynomial--logarithmic transseries \cref{eq:P-ansatz}; recurrence \cref{eq:P-recurrence} computes every coefficient \(\Coeff{\ell^j}{P_r}\).
\item The symmetric translations \(m\mapsto m+1/4\) and \(m\mapsto m+3/4\) produce a common leading constant \(1/\sqrt{2\pi}\), eliminate all odd inverse powers from the logarithmic correction, and yield substantially sharper Lambert carriers.
\item Binet's exact integral supplies a rigorous completion beyond all algebraic orders.  Any terminant or other exponentially improved decomposition is transported to the inverse by the same Lagrange--B\"urmann operator.
\end{enumerate}

Together these results provide a complete, auditable asymptotic framework for the swing factorial: direct values, continuous inverse sheets, discrete threshold inverses, derivatives, complex continuations, exact coefficient generation, and beyond-all-orders refinement all arise from one coherent calculus.

\appendix

\section{Formula sheet}
\label{app:formula-sheet}

For quick reference, the minimal data needed to reproduce the entire calculation are:
\begin{align*}
 \al&=\log4,& K_\eps&=2^\eps/\sqrt\pi,& \be_\eps&=\eps-1/2,\\
 \lambda_r^{(\eps)}
 &=\frac{(-1)^{r+1}}{r(r+1)}
 \left[\B_{r+1}(\eps+1/2)-\B_{r+1}(1)\right],&&&\\
 c_0^{(\eps)}&=1,&
 c_r^{(\eps)}&=\frac1r\sum_{k=1}^rk\lambda_k^{(\eps)}c_{r-k}^{(\eps)},&&\\
 w_0(y)&=-\frac1{2\al}\W_{-1}\!\left(-\frac{2\al}{\pi y^2}\right),&&\\
 w_1(y)&=\frac1{2\al}\W_0\!\left(\frac{\pi\al y^2}{2}\right),&&\\
 d_r^{(\eps)}&=-\frac1\al\Coeff{t^r}{
 \be_\eps\log(1+tD_{<r})+
 H_\eps\!\left(\frac{t}{1+tD_{<r}}\right)},&&&\\
 X_\eps&=\log\!\left(y\al^{\be_\eps}/K_\eps\right),&
 \mu_r^{(\eps)}&=\al^r\lambda_r^{(\eps)}.&&
\end{align*}
The centered data are
\begin{align*}
 s_\eps&=(2\eps+1)/4,&
 \widehat K&=1/\sqrt{2\pi},\\
 \widehat\lambda_{2q-1}^{(\eps)}&=0,&
 \widehat\lambda_{2q}^{(\eps)}&=-\frac{\B_{2q+1}(s_\eps)}{q(2q+1)},\\
 \widehat w_0(y)&=-\frac1{2\al}\W_{-1}\!\left(-\frac\al{\pi y^2}\right),&&\\
 \widehat w_1(y)&=\frac1{2\al}\W_0\!\left(4\pi\al y^2\right).&&
\end{align*}

\section{Additional inverse coefficients}
\label{app:more-d}

For independent implementation checks, the next uncentered coefficients are
\begin{align}
 d_7^{(0)}={}&-
 \frac{510\al^6+539\al^5+357\al^4-525\al^3-4025\al^2+4200\al-840}
 {430080\al^7},
 \label{eq:d7-even}\\
 d_8^{(0)}={}&-
 \frac{6120\al^6+14014\al^5+20454\al^4+15015\al^3-101220\al^2+68040\al-10080}
 {10321920\al^8},
 \label{eq:d8-even}\\
 d_7^{(1)}={}&-\frac{1}{430080\al^7}\Bigl(
 990\al^6+12971\al^5+108549\al^4+224175\al^3
 \nonumber\\[-0.25em]
 &\hspace{34mm}{}+159075\al^2+39480\al+2520\Bigr),
 \label{eq:d7-odd}\\
 d_8^{(1)}={}&\frac{1}{3440640\al^8}\Bigl(
 1680\al^7+59232\al^6+633122\al^5+2123226\al^4
 \nonumber\\[-0.25em]
 &\hspace{24mm}{}+2501205\al^3+1183140\al^2
 +210840\al+10080\Bigr).
 \label{eq:d8-odd}
\end{align}
For the centered normalization,
\begin{align}
 \widehat d_7^{(0)}&=
 \frac{167\al^4-150\al^2+96}{196608\al^6},&
 \widehat d_7^{(1)}&=
 \frac{167\al^4-150\al^2+96}{196608\al^6},
 \label{eq:d7-centered}\\
 \widehat d_8^{(0)}&=-
 \frac{4155\al^6-1840\al^4+1536\al^2-768}{3145728\al^7},&
 \widehat d_8^{(1)}&=
 \frac{4155\al^6-1840\al^4+1536\al^2-768}{3145728\al^7}.
 \label{eq:d8-centered}
\end{align}
Substitution into \cref{eq:fundamental-D}, with the corresponding centered or uncentered \(H\), annihilates every coefficient through the displayed order.

\section{Source and method notes}
\label{app:sources}

The exact sequence normalization and parity identities are those recorded for OEIS A056040.  The shifted gamma expansion, gamma-ratio asymptotics, Binet integral, and Lambert branches follow standard conventions of the NIST Digital Library of Mathematical Functions.  The organizational choice to express exponentiation through Bell polynomials and inversion through coefficient extraction follows the coefficient-calculus and series/transseries drafts in the ProveIt repository.  All special swing-factorial coefficient tables in this article were regenerated from the displayed formulae and checked by exact formal residual substitution; the numerical tables were independently checked with high-precision gamma and Lambert evaluations.

\begin{thebibliography}{99}

\bibitem{OEISA056040}
OEIS Foundation Inc.,
``A056040: Swinging factorial,''
\emph{The On-Line Encyclopedia of Integer Sequences},
\url{https://oeis.org/A056040}, accessed 3 September 2026.

\bibitem{DLMF}
NIST Digital Library of Mathematical Functions,
F.~W.~J. Olver, A.~B. Olde Daalhuis, D.~W. Lozier, B.~I. Schneider,
R.~F. Boisvert, C.~W. Clark, B.~R. Miller, B.~V. Saunders, H.~S. Cohl,
and M.~A. McClain, eds.,
\url{https://dlmf.nist.gov/}; see especially \S\S4.13, 5.9, and 5.11.

\bibitem{Corless1996}
R.~M. Corless, G.~H. Gonnet, D.~E.~G. Hare, D.~J. Jeffrey, and D.~E. Knuth,
``On the Lambert \(W\) function,''
\emph{Advances in Computational Mathematics} \textbf{5} (1996), 329--359.

\bibitem{Gessel2016}
I.~M. Gessel,
``Lagrange inversion,''
\emph{Journal of Combinatorial Theory, Series A} \textbf{144} (2016), 212--249.

\bibitem{Comtet1974}
L.~Comtet,
\emph{Advanced Combinatorics},
D. Reidel, Dordrecht, 1974.

\bibitem{Olver1997}
F.~W.~J. Olver,
\emph{Asymptotics and Special Functions},
AKP Classics, A K Peters, Wellesley, MA, 1997.

\bibitem{ParisKaminski2001}
R.~B. Paris and D.~Kaminski,
\emph{Asymptotics and Mellin--Barnes Integrals},
Cambridge University Press, Cambridge, 2001.

\bibitem{Temme1996}
N.~M. Temme,
\emph{Special Functions: An Introduction to the Classical Functions of Mathematical Physics},
Wiley, New York, 1996.

\bibitem{FlajoletSedgewick2009}
P.~Flajolet and R.~Sedgewick,
\emph{Analytic Combinatorics},
Cambridge University Press, Cambridge, 2009.

\bibitem{ProveItCCC}
V.~Reshetnikov,
``Combinatorial Coefficient Calculus,''
ProveIt repository, semi-formalized research-frontier draft,
\url{https://github.com/VladimirReshetnikov/ProveIt/tree/main/Analysis/FabiusFunction/docs/semi-formalized-research-frontiers/drafts/combinatorial-coefficient-calculus/Combinatorial_Coefficient_Calculus},
accessed 3 September 2026.

\bibitem{ProveItTransseries}
V.~Reshetnikov,
``Series and transseries,''
ProveIt repository, semi-formalized research-frontier drafts,
\url{https://github.com/VladimirReshetnikov/ProveIt/tree/main/Analysis/FabiusFunction/docs/semi-formalized-research-frontiers/drafts/series-and-transseries},
accessed 3 September 2026.

\end{thebibliography}

\end{document}
